\begin{abstract}
This paper presents a geometric framework for analyzing magic squares that bridges discrete combinatorics and continuous mathematics. By embedding the 880 unique $4 \times 4$ magic squares as point clouds in $\mathbb{R}^3$, we reveal that their arithmetic constraints manifest as geometric invariants with physical meaning. We prove a Six-Covariance Theorem: a permutation of $\{1,\ldots,16\}$ forms a magic square if and only if six specific covariances vanish---two first-order conditions capturing \textit{balance} (the center of mass lies at the grid center), two second-order conditions capturing \textit{uniformity} (values are evenly distributed, as measured by moment of inertia), and two diagonal membership conditions ensuring the diagonals are representative samples of the whole. This continuous formulation unifies physical intuition about balance with rigorous statistical characterization. The framework enables classification of squares into 21 distinct convex hull configurations, reveals a highly connected structure under the Earth Mover's Distance metric, and yields a gradient-based optimization algorithm achieving a 73-fold empirical speedup over classical backtracking for $n=4$. Our analysis demonstrates that the ancient ``magic'' of these squares is not mere arithmetic coincidence, but a manifestation of deep geometric balance.
\end{abstract}
